\documentclass{article}
\usepackage{amsmath}
\usepackage{amssymb}
\usepackage{geometry}
\usepackage{booktabs}
\usepackage{graphicx}
\usepackage{xepersian}
\settextfont{XB Niloofar}

\geometry{a4paper, margin=1in}

\title{\lr{Mahalanobis Distance} \\ (فاصله ماهالانوبیس)}
\author{}
\date{}
\begin{figure}[h!]
  \centering
  \includegraphics[width=0.75\textwidth]{img_page_3_1.png}
  \caption{تفاوت هندسی فاصله اقلیدسی و فاصله ماهالانوبیس در توزیع‌های غیرکروی}
\end{figure}

\begin{figure}[h!]
  \centering
  \includegraphics[width=0.75\textwidth]{img_page_4_1.png}
  \caption{تأثیر جهت‌های مختلف واریانس داده‌ها بر محاسبه فاصله}
\end{figure}

\begin{figure}[h!]
  \centering
  \includegraphics[width=0.75\textwidth]{img_page_19_1.png}
  \caption{فواصل نقاط نسبت به مرکز توزیع با در نظر گرفتن همبستگی ویژگی‌ها}
\end{figure}

\begin{figure}[h!]
  \centering
  \includegraphics[width=0.75\textwidth]{img_page_24_1.png}
  \caption{مرز تصمیم‌گیری و دسته‌بندی کلاس‌ها بر اساس فاصله ماهالانوبیس}
\end{figure}










\begin{document}

\maketitle

\section*{فهرست مطالب}
\begin{enumerate}
    \item چرا به فاصله ماهالانوبیس نیاز داریم؟
    \item تاثیر مقیاس‌بندی ویژگی‌ها
    \item مقدمه‌ای بر ماتریس کوواریانس (\lr{Covariance Matrix})
    \item تفسیر هندسی ماتریس کوواریانس
    \item تبدیل‌های خطی و کوواریانس
    \item تبدیل سفیدسازی (\lr{Whitening Transformation})
    \item فاصله ماهالانوبیس
    \item ماتریس کوواریانس در تئوری تصمیم‌گیری بیز (\lr{Bayesian Decision Theory})
    \item مقایسه: \lr{PCA} در مقابل سفیدسازی
    \item نتیجه‌گیری و نکات کلیدی
\end{enumerate}

\newpage

\section*{۱. چرا به فاصله ماهالانوبیس نیاز داریم؟}
در بسیاری از کاربردهای دنیای واقعی، ما با داده‌های چندمتغیره مواجه هستیم که در آنها ویژگی‌ها اغلب همبسته هستند و مقیاس‌های متفاوتی دارند. در حالی که فاصله اقلیدسی (\lr{Euclidean distance}) یک انتخاب متداول برای اندازه‌گیری فاصله‌ها است، یک محدودیت بزرگ دارد: فرض می‌کند که همه ویژگی‌ها غیرهمبسته بوده و دارای واریانس یکسانی هستند. این فرض به ندرت در عمل برقرار است.

مجموعه‌داده‌ای را در نظر بگیرید که یک ویژگی آن به متر و ویژگی دیگر به کیلومتر اندازه‌گیری می‌شود. فاصله اقلیدسی به طور نامتناسبی وزن بیشتری به ویژگی با مقادیر بزرگتر می‌دهد. به همین ترتیب، اگر ویژگی‌ها همبسته باشند، شکل توزیع داده‌ها به جای کره، یک بیضی‌گون است. فاصله اقلیدسی در به حساب آوردن این ساختار زمینه‌ای شکست می‌خورد که منجر به نتایج گمراه‌کننده در دسته‌بندی و خوشه‌بندی می‌شود.

فاصله ماهالانوبیس (\lr{Mahalanobis distance}) این مشکل را با اندازه‌گیری فاصله بین یک نقطه و یک توزیع حل می‌کند، که توسط واریانس هر ویژگی و همبستگی بین آنها نرمال‌سازی می‌شود.

\subsection*{مثال تشریحی: فاصله اقلیدسی در مقابل ماهالانوبیس}
نقطه ۱ و نقطه ۲ فاصله اقلیدسی یکسانی از مرکز دارند، اما در ``نواحی احتمال'' متفاوتی قرار می‌گیرند. نقطه ۲ در جهت واریانس کمتر (عمود بر همبستگی اصلی) قرار دارد که آن را از نظر آماری غیرمحتمل‌تر (دورتر در فاصله ماهالانوبیس) از نقطه ۱ می‌کند.

\section*{۲. تاثیر مقیاس‌بندی ویژگی‌ها}
\subsection*{مطالعه موردی: شباهت موشک‌های ایرانی}
در این مثال، از مجموعه‌داده موشک‌های بالستیک ایرانی برای حل یک مسئله شباهت استفاده می‌کنیم. می‌خواهیم تعیین کنیم کدام موشک بر اساس برد و سرعت، شباهت بیشتری به خیبر شکن (موشک A) دارد.

مجموعه‌داده:
\begin{itemize}
    \item \textbf{موشک A (خیبر شکن):} برد ۱۴۵۰ کیلومتر، سرعت ۵ ماخ (هدف)
    \item \textbf{موشک B (سجیل):} برد ۲۰۰۰ کیلومتر، سرعت ۱۰ ماخ
    \item \textbf{موشک C (فتاح):} برد ۱۴۰۰ کیلومتر، سرعت ۱۴ ماخ
\end{itemize}

\subsubsection*{رویکرد ساده‌لوحانه (فاصله خام اقلیدسی)}
فاصله اقلیدسی استاندارد بین هدف (خیبر شکن) و سایر موشک‌ها به صورت زیر است:
\begin{itemize}
    \item خیبر در مقابل سجیل:
    \[
    D = \sqrt{550^2 + 5^2} \approx 550.02
    \]
    \item خیبر در مقابل فتاح:
    \[
    D = \sqrt{50^2 + 9^2} \approx 50.80
    \]
\end{itemize}
\textbf{مشکل:} موشک فتاح بسیار نزدیک‌تر (۵۰.۸) به نظر می‌رسد تا سجیل (۵۵۰.۰). اما در واقعیت، ۵۰ کیلومتر تفاوت در برد موشک ناچیز است، اما ۹ ماخ تفاوت در کلاس سرعت (هایپرسونیک در مقابل متوسط) بسیار بزرگ است. محاسبات خام توسط ویژگی برد تحت‌الشعاع قرار می‌گیرد.

\subsubsection*{رویکرد تصحیح‌شده (مقیاس‌بندی استاندارد)}
برای رفع این مشکل، داده‌ها را با استفاده از امتیاز استاندارد (\lr{Z-scores}) نرمال‌سازی می‌کنیم. این کار هر دو ویژگی را در یک مقیاس قرار می‌دهد. پس از مقیاس‌بندی، فواصل به صورت زیر در می‌آیند:
\begin{itemize}
    \item خیبر در مقابل سجیل: $D \approx 2.44$
    \item خیبر در مقابل فتاح: $D \approx 2.45$
\end{itemize}
\textbf{نتیجه:} اکنون سجیل به درستی به عنوان نزدیک‌ترین گزینه به خیبر شناسایی می‌شود. مقیاس‌بندی، تفاوت سرعت بزرگ فتاح را جریمه کرده و ویژگی‌ها را متعادل می‌کند.

\textbf{نکته کلیدی:} عدم تطابق واحدها یک مشکل خاموش در یادگیری ماشین است. همیشه ویژگی‌های خود را قبل از محاسبه فواصل، آموزش شبکه‌های عصبی یا خوشه‌بندی مقیاس‌بندی کنید.

\section*{۳. مقدمه‌ای بر ماتریس کوواریانس}
ماتریس کوواریانس (\lr{covariance matrix}) یک ماتریس مربع است که واریانس‌ها و کوواریانس‌های مجموعه‌ای از متغیرهای تصادفی را خلاصه می‌کند. برای داده‌های با $d$ ویژگی، ماتریس کوواریانس $\Sigma$ یک ماتریس $d \times d$ به صورت زیر است:
\[
\Sigma = \begin{bmatrix} 
\sigma_{11} & \sigma_{12} & \cdots & \sigma_{1d} \\
\sigma_{21} & \sigma_{22} & \cdots & \sigma_{2d} \\
\vdots & \vdots & \ddots & \vdots \\
\sigma_{d1} & \sigma_{d2} & \cdots & \sigma_{dd} 
\end{bmatrix}
\]
که در آن:
\begin{itemize}
    \item $\sigma_{ii} = Var(x_i)$: واریانس ویژگی $i$-ام (عناصر روی قطر اصلی).
    \item $\sigma_{ij} = Cov(x_i, x_j)$: کوواریانس بین ویژگی‌های $i$-ام و $j$-ام (عناصر خارج از قطر).
\end{itemize}
\begin{figure}[h!]
  \centering
  \includegraphics[width=0.75\textwidth]{img_page_9_1.png}
  \caption{پراکندگی داده‌ها و نسبت فواصل ماهالانوبیس نقاط مختلف}
\end{figure}




ویژگی‌ها: متقارن است ($\Sigma = \Sigma^T$) و نیمه‌معین مثبت است (تمام مقادیر ویژه غیرمنفی هستند $\lambda_i \ge 0$).

\section*{۴. تفسیر هندسی ماتریس کوواریانس}
ماتریس کوواریانس شکل و جهت توزیع داده‌ها را در فضای ویژگی تعریف می‌کند:
\begin{enumerate}
    \item \textbf{بردارهای ویژه (\lr{Eigenvectors}):} بردارهای ویژه $\Sigma$ نشان‌دهنده محورهای اصلی توزیع داده‌ها هستند و جهت‌های حداکثر واریانس را مشخص می‌کنند.
    \item \textbf{مقادیر ویژه (\lr{Eigenvalues}):} مقادیر ویژه متناظر، مقدار واریانس را در امتداد جهت‌های این بردارهای ویژه نشان می‌دهند.
    \item \textbf{بیضی‌گون:} در ۲ بعد، ماتریس کوواریانس یک بیضی و در $d$ بعد، یک بیضی‌گون با ابعاد بالا (\lr{hyper-ellipsoid}) را تعریف می‌کند. طول محورهای این بیضی‌گون متناسب با جذر مقادیر ویژه است.
\end{enumerate}

\section*{۵. تبدیل‌های خطی و کوواریانس}
کوواریانس تحت تبدیل‌های خطی دستخوش تغییر می‌شود.

\textbf{قضیه ۱ (ویژگی تبدیل کوواریانس):} 
فرض کنید $x \in \mathbb{R}^d$ یک بردار تصادفی با میانگین $\mu$ و ماتریس کوواریانس $\Sigma$ باشد.
فرض کنید $W \in \mathbb{R}^{d \times d}$ یک ماتریس تبدیل ثابت باشد. اگر بردار تصادفی تبدیل‌شده $y$ به صورت $y = W^T(x - \mu)$ تعریف شود، آنگاه ماتریس کوواریانس بردار تبدیل‌شده به صورت زیر خواهد بود:
\[
Cov(y) = W^T \Sigma W
\]

\textbf{اثبات:}
ابتدا میانگین $y$ را محاسبه می‌کنیم:
\[
E[y] = E[W^T(x-\mu)] = W^T(E[x]-\mu) = W^T(\mu-\mu) = 0
\]
سپس کوواریانس $y$:
\[
Cov(y) = E[(y - E[y])(y - E[y])^T] = E[y y^T] = E[W^T(x-\mu)(x-\mu)^T W]
\]
با استفاده از خطی بودن عملگر امید ریاضی، ماتریس‌های ثابت را بیرون می‌آوریم:
\[
Cov(y) = W^T E[(x-\mu)(x-\mu)^T] W = W^T \Sigma W
\]
این قضیه پایه و اساس سفیدسازی داده‌ها است. با انتخاب $W$ به طوری که $W^T \Sigma W = I$ شود، داده‌ها غیرهمبسته و دارای واریانس واحد می‌شوند.

\section*{۶. تبدیل سفیدسازی (\lr{Whitening Transformation})}
سفیدسازی داده‌ها را به گونه‌ای تبدیل می‌کند که ماتریس کوواریانس آن به ماتریس همانی ($I$) تبدیل شود. 
\begin{figure}[h!]
  \centering
  \includegraphics[width=0.75\textwidth]{img_page_8_1.png}
  \caption{بیضی‌های هم‌فاصله ماهالانوبیس بر اساس بردارهای ویژه ماتریس کوواریانس}
\end{figure}




ما به دنبال ماتریس تبدیل $W$ هستیم به طوری که $Cov(y) = W^T \Sigma W = I$ شود.
با استفاده از تجزیه مقادیر ویژه ماتریس کوواریانس داریم: $\Sigma = \Phi \Lambda \Phi^T$ که در آن:
\begin{itemize}
    \item $\Phi$: ماتریس متعامد بردارهای ویژه.
    \item $\Lambda$: ماتریس قطری مقادیر ویژه.
\end{itemize}
ما $W = \Phi \Lambda^{-1/2}$ را انتخاب می‌کنیم. راستی‌آزمایی:
\[
Cov(y) = (\Phi \Lambda^{-1/2})^T \Sigma (\Phi \Lambda^{-1/2}) = \Lambda^{-1/2} \Phi^T (\Phi \Lambda \Phi^T) \Phi \Lambda^{-1/2} = I
\]

\textbf{تفسیر هندسی:}
\begin{enumerate}
    \item \textbf{چرخش ($\Phi^T$):} داده‌ها را با محورهای اصلی تراز کرده و همبستگی‌ها را حذف می‌کند.
    \item \textbf{مقیاس‌بندی ($\Lambda^{-1/2}$):} داده‌ها را در امتداد این محورها می‌کشد یا فشرده می‌کند تا طول همه محورها ۱ شود.
\end{enumerate}
در این فضای سفیدشده، فاصله اقلیدسی معادل با فاصله ماهالانوبیس در فضای اصلی است.

\section*{۷. فاصله ماهالانوبیس}
فاصله ماهالانوبیس بین یک نقطه $x$ و یک توزیع با میانگین $\mu$ و ماتریس کوواریانس $\Sigma$ به صورت زیر تعریف می‌شود:
\[
D_M(x, \mu) = \sqrt{(x - \mu)^T \Sigma^{-1} (x - \mu)}
\]

\textbf{اثبات هم‌ارزی:}
فرض کنید $y = W^T(x-\mu)$ که در آن $W$ ماتریس سفیدکننده است. فاصله اقلیدسی در فضای سفیدشده برابر است با:
\[
\|y\|_2 = \sqrt{y^T y} = \sqrt{(x-\mu)^T W W^T (x-\mu)}
\]
نشان می‌دهیم که $W W^T = \Sigma^{-1}$ است. داریم $W = \Phi \Lambda^{-1/2}$، پس:
\[
W W^T = \Phi \Lambda^{-1/2} (\Phi \Lambda^{-1/2})^T = \Phi \Lambda^{-1} \Phi^T
\]
از طرفی، معکوس $\Sigma = \Phi \Lambda \Phi^T$ با توجه به متعامد بودن $\Phi$ برابر است با:
\[
\Sigma^{-1} = \Phi \Lambda^{-1} \Phi^T
\]
بنابراین $W W^T = \Sigma^{-1}$ برقرار است و فاصله اقلیدسی در فضای سفیدشده با فاصله ماهالانوبیس در فضای اصلی برابر است.

\section*{۸. ماتریس کوواریانس در تئوری تصمیم‌گیری بیز}
توزیع گاوسی چندمتغیره به صورت زیر است:
\[
p(x|\omega_i) = \frac{1}{(2\pi)^{d/2} |\Sigma_i|^{1/2}} \exp \left( -\frac{1}{2} (x - \mu_i)^T \Sigma_i^{-1} (x - \mu_i) \right)
\]
جمله داخل نما دقیقاً نصف فاصله ماهالانوبیس مربع است. معکوس ماتریس کوواریانس $\Sigma^{-1}$ ویژگی‌ها را با واریانس آنها مقیاس‌بندی کرده و فضا را برای رفع همبستگی می‌چرخاند.

\section*{آیا سفیدسازی با نرمال‌سازی تک‌ویژگی معادل است؟}
خیر. نرمال‌سازی تک‌ویژگی (استانداردسازی) هر ویژگی را به طور مستقل به واریانس واحد می‌رساند اما همبستگی بین آنها را حفظ می‌کند. اما سفیدسازی، همبستگی‌ها را نیز به طور کامل حذف می‌کند.

\section*{۹. مقایسه: \lr{PCA} در مقابل سفیدسازی}
\begin{itemize}
    \item \textbf{\lr{PCA} (فقط چرخش):} داده‌ها را روی محورهای اصلی تصویر می‌کند. همبستگی‌ها را حذف می‌کند اما مقیاس واریانس را حفظ می‌کند. ماتریس کوواریانس حاصل قطری ($\Lambda$) است.
    \item \textbf{سفیدسازی (چرخش + مقیاس‌بندی):} چرخش PCA را اعمال کرده و سپس واریانس هر مولفه را به ۱ می‌رساند. ماتریس کوواریانس حاصل همانی ($I$) است.
\end{itemize}

\begin{table}[h]
\centering
\caption{تفاوت‌های کلیدی بین \lr{PCA} و سفیدسازی}
\vspace{1mm}
\begin{tabular}{lll}
\toprule
\textbf{ویژگی} & \textbf{\lr{PCA}} & \textbf{سفیدسازی} \\
\midrule
هدف اصلی & کاهش ابعاد / بیشینه‌سازی واریانس & نرمال‌سازی واریانس و همبستگی \\
نوع تبدیل & فقط چرخش & چرخش + مقیاس‌بندی \\
واریانس خروجی & مقادیر ویژه ($\lambda_i$) & واریانس واحد ($1$) \\
ماتریس کوواریانس & قطری ($\Lambda$) & همانی ($I$) \\
شکل هندسی & بیضی هم‌جهت با محورها & کره کامل (همسانگرد) \\
\bottomrule
\end{tabular}
\end{table}

\section*{۱۰. نتیجه‌گیری و نکات کلیدی}
\begin{itemize}
    \item فاصله اقلیدسی داده‌ها را کروی فرض می‌کند، در حالی که فاصله ماهالانوبیس داده‌ها را بیضوی فرض می‌کند.
    \item ماتریس کوواریانس شکل، مقیاس و جهت توزیع داده‌ها را ثبت می‌کند.
    \item سفیدسازی داده‌ها را کروی می‌کند به طوری که فاصله اقلیدسی در فضای سفیدشده با فاصله ماهالانوبیس در فضای اصلی برابر می‌شود.
    \item این فاصله برای تشخیص آنومالی و تعیین مرزهای تصمیم‌گیری بیز ضروری است.
\end{itemize}

\end{document}
